% ----------------------------------------------------------------------
% Why:     Role-targeted resume variant written against the Rockstar Games
%          "Online Systems Administrator" req (Bengaluru, 24x7 NOC for the
%          online games platform). Every other variant sells feature or AI
%          delivery. This req is a run-the-platform role: keep a public-facing
%          service up, monitor it, find the root cause, write the runbook, and
%          automate the repetitive parts over REST. Those things are genuinely
%          in the repositories and are invisible in every other variant.
% What:    Source for public/resume-rockstar-sysadmin.pdf. Per CLAUDE.md
%          section 6 this is a role-targeted variant: it drifts from
%          resume.tex on purpose and must NOT be "synced" back to it.
% Result:  One page. Every claim traces to docs/resume-evidence.json.
%          DELIBERATELY SILENT on the half of the req that no repository on
%          disk supports (section 1.3): Windows Server administration, Active
%          Directory, SQL Server 2017/2019 clustering, VMware, IIS/.NET C#,
%          PowerShell/DSC, Chef, Check MK, Splunk, Grafana and Cisco IOS.
%          Claiming any of them would be an invented keyword. The covering
%          letter is where that gap gets addressed honestly, not the resume.
% Changed: 2026-10-05 - created for the Rockstar Games Online Systems
%          Administrator req.
% Build:   bash ~/.claude/skills/resume-forge/scripts/build_resume.sh \
%            public/resume-rockstar-sysadmin.tex
% ----------------------------------------------------------------------
\documentclass[10pt,a4paper]{article}

% ----------------------------------------------------------------------
% Packages
% ----------------------------------------------------------------------
\usepackage[T1]{fontenc}
\usepackage[utf8]{inputenc}
\usepackage{lmodern}
\usepackage[a4paper,top=0.9cm,bottom=0.8cm,left=1.25cm,right=1.25cm]{geometry}
\usepackage{titlesec}
\usepackage{enumitem}
\usepackage[hidelinks]{hyperref}

\hypersetup{
  pdftitle={Harish G - Online Systems Administrator Resume},
  pdfauthor={Harish G},
  colorlinks=false
}

% ----------------------------------------------------------------------
% Section styling: sans-serif bold heading with a thin rule beneath
% ----------------------------------------------------------------------
\titleformat{\section}
  {\normalsize\sffamily\bfseries}{}{0em}{}[\titlerule]
\titlespacing*{\section}{0pt}{8pt}{3pt}

% Compact bullet lists: the single biggest lever on whether this fits a page.
\setlist[itemize]{leftmargin=1.2em,itemsep=1.5pt,topsep=2pt,parsep=0pt}

\setlength{\parindent}{0pt}
\linespread{0.97}\selectfont
\pagestyle{empty}

% ----------------------------------------------------------------------
% Helpers
% ----------------------------------------------------------------------
\newcommand{\entry}[2]{\textbf{#1}\hfill{#2}\par}
\newcommand{\sep}{\,\textperiodcentered\,}
\newlength{\skillwidth}\setlength{\skillwidth}{6.4em}
\newlength{\skilllabel}
\newcommand{\skill}[2]{%
  \settowidth{\skilllabel}{\textbf{#1}}%
  \hangindent=\skillwidth\hangafter=1%
  \mbox{\textbf{#1}}%
  \ifdim\skilllabel<\skillwidth\hspace{\dimexpr\skillwidth-\skilllabel}\else\hspace{0.6em}\fi
  #2\par}

\begin{document}

% ----------------------------------------------------------------------
% Header
% ----------------------------------------------------------------------
{\LARGE\bfseries Harish G}\\[1pt]
{\normalsize\bfseries Systems \& Platform Operations \textbar{} Production Support, Monitoring \& Automation}\\[2pt]
{\small Bengaluru, India \quad\textbar\quad +91 7892855850 \quad\textbar\quad
\href{mailto:harishgreddy.work@gmail.com}{harishgreddy.work@gmail.com} \quad\textbar\quad
\href{https://www.linkedin.com/in/harishgreddy/}{LinkedIn} \quad\textbar\quad
\href{https://github.com/harish00506/}{GitHub} \quad\textbar\quad \href{https://www.harishg.com}{harishg.com}}

% ----------------------------------------------------------------------
% Profile
% ----------------------------------------------------------------------
\section{Profile}
I stay with a system after it ships. On a public-facing telephony platform I own the call pipeline
in production: fixed outbound routing, added reconnect handling for dropped Gemini 1006 sockets, and
traced a Docker TLS failure back to NAT hairpinning on a self-hosted domain, which read as an
application bug until somebody followed the packet path. I build the observability in rather than bolt
it on, and write the runbook that says what to do at 2am. Software Developer at CortexCraft.ai.

% ----------------------------------------------------------------------
% Core Skills
% ----------------------------------------------------------------------
\section{Core Skills}
\skill{Operations}{24x7 ownership of a live, public-facing call pipeline \sep root-cause analysis
from symptom to packet path \sep runbooks, a hotfix flow and a rollback procedure \sep backend and
provider migrations run without interrupting live traffic}
\skill{Monitoring}{single-line structured JSON logging \sep \texttt{X-Request-ID} correlated end to
end \sep per-response status and duration in milliseconds \sep Prometheus instrumentator on
\texttt{/metrics} \sep health endpoints \sep a real-time WebSocket monitoring
dashboard over a Windows desktop agent \sep append-only audit logs with TTL retention}
\skill{Platform}{Linux hosts under Docker and Docker Compose \sep Jenkins, GitLab CI and GitHub
Actions as the merge gate \sep containerised multi-service deploys \sep TLS, DNS and NAT fault
isolation \sep configuration in the database, so a change is a config row and not a deploy}
\skill{Scripting}{Python \sep Java \sep JavaScript / TypeScript \sep SQL \sep Celery workers and
scheduled jobs \sep bulk CSV and Excel import / export jobs}
\skill{Databases}{PostgreSQL with Alembic and Flyway versioned migrations \sep MongoDB \sep Redis
(locks and queues) \sep Neo4j \sep credentials encrypted at rest}
\skill{Automation}{External REST API that other products drive \sep webhook callback handlers \sep
ticket-to-pull-request workflows over the Jira and GitHub APIs \sep 132 pytest files acting as the
CI merge gate}

% ----------------------------------------------------------------------
% Experience
% ----------------------------------------------------------------------
\section{Experience}

\entry{Software Developer, CortexCraft.ai}{Jan 2026 -- Present}
\begin{itemize}
  \item Authored 589 of 626 commits on the voice-agent SaaS platform and 67 of 75 on its core
        engine, and own production reliability of the call pipeline it runs: fixed outbound routing,
        added reconnect handling for dropped Gemini 1006 sockets, and traced a Docker TLS failure to
        NAT hairpinning on a self-hosted domain.
  \item Lead author (172 of 192 commits) on the multilingual survey platform. Migrated it off a Node
        backend onto FastAPI and off one speech provider onto another behind a config-driven
        registry of 24 locales, both without interrupting surveys that were running at the time.
  \item Run the build and deploy path behind it: Jenkins and GitLab CI, Docker Compose, 132 pytest
        files gating every merge, and Alembic-migrated PostgreSQL treated as a release.
  \item Maintain integration workflows against two telephony vendors and several partner APIs: an
        external REST API other products call, webhook callback handlers, per-tenant credentials
        encrypted at rest, and an agent that turns Jira tickets into reviewed pull requests.
\end{itemize}

\vspace{1pt}
\entry{Freelance Web \& App Developer (Self-Employed)}{2023 -- 2026}
Delivered client systems end to end, then supported them: deployment, environment changes, upgrades
and fixes, with direct client contact for triage and release sign-off.

% ----------------------------------------------------------------------
% Projects
% ----------------------------------------------------------------------
\section{Selected Projects}

\textbf{CortexCraft Voice Agent, production telephony platform} \textit{Python, FastAPI, WebSocket, Docker, React}
\begin{itemize}
  \item 589 of 626 commits. One platform behind several shipped products, so a fix to the call
        pipeline reaches all of them at once. Most of the real work was keeping it up: routing,
        socket reconnects and TLS, each traced to a cause rather than restarted past.
\end{itemize}

\vspace{1pt}
\textbf{KisanVoice AI, multilingual survey platform} \textit{Python, FastAPI, MongoDB, Jenkins, Docker Compose}
\begin{itemize}
  \item 172 of 192 commits. Operations built in rather than added later: every request carries an
        \texttt{X-Request-ID} that is forwarded or generated and echoed back, every log line is a
        single JSON object carrying that ID, each response is logged with its status and duration in
        milliseconds, and a Prometheus instrumentator exposes \texttt{/metrics}.
  \item Ships with a deployment runbook covering prerequisites, environment setup, a verification
        step and a symptom-to-fix section for the failures the team actually hit, plus a written
        release process with a hotfix flow and a rollback procedure. 24 configured locales come from
        a config-driven registry, so adding one is a config row rather than a deploy.
\end{itemize}

\vspace{1pt}
\textbf{LeadCall AI, multi-tenant outbound calling} \textit{FastAPI, Celery, PostgreSQL, Redis, React}
\begin{itemize}
  \item Queued the slow work instead of blocking on it: Celery dial and callback workers behind
        Redis dial locks, per-tenant credentials encrypted at rest, bulk Excel import and call-log
        export, and an append-only audit log of every operator action.
\end{itemize}

% ----------------------------------------------------------------------
% Education
% ----------------------------------------------------------------------
\section{Education \& Ways of Working}
\textbf{B.E., Information Science and Engineering}, Vivekananda Institute of Technology, Bangalore.
One issue per branch, tests as the merge gate, a runbook entry for every failure seen twice.

\end{document}
